\section{Machine learning for post-EVAR outcomes}
\begin{itemize}
    \item Anatomical and radiomic features for endoleak prediction \cite{peng_predict_2024}.
    \item Point-cloud network with morphology, radiomics and CFD (AUC 0.81 to 0.93). \citeneeded{Peng et al.\ 2026, medRxiv}
\end{itemize}

\section{CFD and machine learning in AAA}
\begin{itemize}
    \item Growth prediction from hemodynamic and multiphysical features \cite{forneris_ai-powered_2023, kim_deep_2023}.
    \item Spatial patterns of hemodynamic indices \cite{rezaeitaleshmahalleh_characterization_2023}.
    \item Deep learning surrogates of CFD \cite{faisal_rapid_2025, fedotova_automatically_2021}.
\end{itemize}

\section{Gap}
\begin{itemize}
    \item In Peng et al.\ the CFD gain is confounded with architecture changes (1 vs.\ 3 feature extractors). \citeneeded{Peng et al.\ 2026}
    \item One test set, no confidence intervals or paired tests.
    \item No study isolates the incremental value of preoperative CFD for post-EVAR complications with a fixed model.
\end{itemize}
This thesis fills the gap with fixed-model comparisons, explicit controls and quantified uncertainty.
